\documentclass[12pt,a4paper]{article}
\usepackage[UTF8]{ctex}
\usepackage{amsmath,amssymb,amsthm}
\usepackage{geometry}
\usepackage{bm}
\usepackage{hyperref}
\usepackage{tikz}
\usetikzlibrary{arrows,arrows.meta,positioning,calc,shapes.geometric}

\geometry{margin=2.5cm}

\title{12.1.4 物体集合的详细理解}
\author{第12章抓取与操作补充说明}
\date{}

\newtheorem{theorem}{定理}
\newtheorem{definition}{定义}
\newtheorem{example}{例子}

\begin{document}

\maketitle

\section{章节概述}

12.1.4章节"物体集合"（Collections of Bodies）将前面的讨论从\textbf{单个物体}扩展到\textbf{多个物体}的情况。这是接触运动学的一个重要推广，因为在实际应用中，我们经常需要处理多个物体同时接触的情况。

\section{从单个物体到多个物体}

\subsection{12.1.3章节的回顾}

在12.1.3章节中，我们讨论了：
\begin{itemize}
\item \textbf{单个物体}$A$受到多个接触的约束
\item 可行twist集合：$\mathcal{V} = \{V_A | F_i^T (V_A - V_{j(i)}) \geq 0 \text{ 对所有 } i\}$
\item 可行集合是$V_A$空间（6维）中的多面体
\end{itemize}

\subsection{12.1.4章节的扩展}

在12.1.4章节中，我们考虑：
\begin{itemize}
\item \textbf{多个物体}同时存在，它们之间相互接触
\item 每个物体都有自己的twist $V_i$
\item 需要同时考虑所有物体的约束
\end{itemize}

\section{两个物体接触的情况}

\subsection{基本设置}

考虑两个物体$i$和$j$在点$p$处接触：
\begin{itemize}
\item 接触法线$\hat{n}$指向物体$i$
\item 接触wrench：$F = ([p]\hat{n}, \hat{n})$
\item 物体$i$的twist：$V_i$
\item 物体$j$的twist：$V_j$
\end{itemize}

\subsection{不可穿透约束}

两个物体之间的不可穿透约束为：
\begin{equation}
F^T (V_i - V_j) \geq 0
\label{eq:two_bodies_constraint}
\end{equation}

这与12.1.2节中的公式(12.5)相同，但现在我们明确考虑两个物体都在运动。

\subsection{复合twist空间}

\textbf{关键概念}：我们需要在\textbf{复合twist空间}中考虑约束。

对于两个物体：
\begin{itemize}
\item 物体$i$的twist：$V_i \in \mathbb{R}^6$（6维）
\item 物体$j$的twist：$V_j \in \mathbb{R}^6$（6维）
\item \textbf{复合twist}：$(V_i, V_j) \in \mathbb{R}^{12}$（12维）
\end{itemize}

约束$F^T (V_i - V_j) \geq 0$在12维复合twist空间中定义了一个\textbf{齐次半空间约束}。

\subsection{为什么是齐次约束？}

约束$F^T (V_i - V_j) \geq 0$可以写成：
\begin{equation}
F^T V_i - F^T V_j \geq 0
\end{equation}

或者：
\begin{equation}
[F^T \quad -F^T] \begin{bmatrix} V_i \\ V_j \end{bmatrix} \geq 0
\end{equation}

这是一个\textbf{齐次约束}，因为：
\begin{itemize}
\item 约束超平面通过原点：$F^T V_i - F^T V_j = 0$在$(V_i, V_j) = (0, 0)$处成立
\item 如果$(V_i, V_j)$满足约束，那么$\lambda (V_i, V_j)$（$\lambda \geq 0$）也满足约束
\end{itemize}

\section{多个物体装配的情况}

\subsection{基本设置}

考虑一个由$m$个物体组成的装配（assembly）：
\begin{itemize}
\item 物体编号：$i = 1, 2, \ldots, m$
\item 每个物体$i$的twist：$V_i \in \mathbb{R}^6$
\item 物体之间可能有多个接触
\end{itemize}

\subsection{复合twist空间}

对于$m$个物体：
\begin{itemize}
\item \textbf{复合twist空间}：$(V_1, V_2, \ldots, V_m) \in \mathbb{R}^{6m}$
\item 对于平面物体：$(V_1, V_2, \ldots, V_m) \in \mathbb{R}^{3m}$（每个物体3维）
\end{itemize}

\textbf{维度分析}：
\begin{itemize}
\item 每个物体有6个自由度（3D）或3个自由度（平面）
\item $m$个物体总共有$6m$个自由度（3D）或$3m$个自由度（平面）
\item 复合twist空间就是所有可能运动的集合
\end{itemize}

\subsection{成对接触约束}

如果物体$i$和物体$j$在点$p$处接触：
\begin{itemize}
\item 接触法线$\hat{n}$指向物体$i$
\item 接触wrench：$F = ([p]\hat{n}, \hat{n})$
\item 约束：$F^T (V_i - V_j) \geq 0$
\end{itemize}

这个约束在$6m$维复合twist空间中定义了一个半空间。

\subsection{所有约束的组合}

假设装配中有多个成对接触，每个接触贡献一个约束。所有约束的\textbf{交集}定义了可行twist集合：

\begin{equation}
\mathcal{V} = \{(V_1, \ldots, V_m) | \text{所有接触约束都满足}\}
\label{eq:assembly_feasible}
\end{equation}

\subsection{可行twist集合的性质}

如果所有物体都是\textbf{不受控制的}（uncontrolled），即没有外部运动控制：
\begin{itemize}
\item 所有约束都是齐次的
\item 可行twist集合是一个\textbf{以原点为根部的多面体凸锥}
\item 原点$(V_1, \ldots, V_m) = (0, \ldots, 0)$总是可行的（所有物体静止）
\end{itemize}

\section{受控物体与非受控物体}

\subsection{基本概念}

在实际应用中，有些物体的运动是\textbf{受控的}（controlled），有些是\textbf{非受控的}（uncontrolled）：

\begin{itemize}
\item \textbf{受控物体}：运动由外部控制（例如，机器人手指）
  \begin{itemize}
  \item 它们的twist $V_i$是\textbf{已知的}或\textbf{指定的}
  \item 例如：$V_{\text{finger}} = [0, 0, 0, v_x, 0, 0]^T$（手指以速度$v_x$移动）
  \end{itemize}

\item \textbf{非受控物体}：运动由接触约束和动力学决定（例如，被抓取的物体）
  \begin{itemize}
  \item 它们的twist $V_i$是\textbf{未知的}，需要求解
  \item 约束限制了它们的可能运动
  \end{itemize}
\end{itemize}

\subsection{约束的变化}

当有受控物体时，约束的性质发生变化：

\begin{itemize}
\item \textbf{齐次约束}：$F^T (V_i - V_j) \geq 0$
  \begin{itemize}
  \item 如果$V_j$是受控的且$V_j \neq 0$，约束变为：$F^T V_i \geq F^T V_j$
  \item 这是一个\textbf{非齐次约束}，因为超平面不通过原点
  \end{itemize}

\item \textbf{可行集合的变化}：
  \begin{itemize}
  \item 如果所有物体都非受控：可行集合是\textbf{以原点为根部的锥}
  \item 如果有受控物体：可行集合是\textbf{一般的多面体}（可能是有界的）
  \end{itemize}
\end{itemize}

\subsection{例子：单个受控物体}

考虑两个物体：
\begin{itemize}
\item 物体1：受控（机器人手指），$V_1$已知
\item 物体2：非受控（被抓取物体），$V_2$未知
\end{itemize}

约束：$F^T (V_1 - V_2) \geq 0$

可以写成：
\begin{equation}
F^T V_2 \leq F^T V_1
\end{equation}

这是一个\textbf{非齐次约束}，因为：
\begin{itemize}
\item 超平面$F^T V_2 = F^T V_1$不通过原点（除非$V_1 = 0$）
\item 可行集合不再是锥，而是一个半空间
\end{itemize}

\section{接触模式}

\subsection{装配的接触模式}

对于整个装配，\textbf{接触模式}（contact mode）是：
\begin{itemize}
\item 装配中每个接触处接触标签的连接
\item 例如：如果有3个接触，接触模式可能是"RRS"、"BSB"等
\end{itemize}

\subsection{与单个物体的区别}

\begin{itemize}
\item \textbf{单个物体}：接触模式描述物体$A$与多个其他物体的接触状态
\item \textbf{多个物体}：接触模式描述整个装配中所有接触的状态
\end{itemize}

\section{具体例子}

\subsection{例子1：两个静止物体接触}

考虑两个物体$A$和$B$在点$p$处接触，且都静止：

\begin{itemize}
\item 物体$A$的twist：$V_A$
\item 物体$B$的twist：$V_B$
\item 约束：$F^T (V_A - V_B) \geq 0$
\end{itemize}

\textbf{可行twist集合}：
\begin{itemize}
\item 复合twist空间：$(V_A, V_B) \in \mathbb{R}^{12}$
\item 约束在12维空间中定义了一个半空间
\item 可行集合是一个以原点为根部的锥
\item 原点$(V_A, V_B) = (0, 0)$是可行的（两个物体都静止）
\end{itemize}

\subsection{例子2：机器人手指抓取物体}

考虑：
\begin{itemize}
\item 物体1：机器人手指（受控），$V_1 = [0, 0, 0, v_x, 0, 0]^T$（以速度$v_x$移动）
\item 物体2：被抓取物体（非受控），$V_2$未知
\item 接触约束：$F^T (V_1 - V_2) \geq 0$
\end{itemize}

\textbf{约束分析}：
\begin{itemize}
\item 约束可以写成：$F^T V_2 \leq F^T V_1$
\item 这是一个非齐次约束
\item 可行集合是$V_2$空间中的一个半空间（不是锥）
\item 如果$v_x > 0$，可行集合不包含原点
\end{itemize}

\subsection{例子3：三个物体的装配}

考虑三个物体$A$、$B$、$C$：
\begin{itemize}
\item $A$和$B$在点$p_1$处接触
\item $B$和$C$在点$p_2$处接触
\item 所有物体都非受控
\end{itemize}

\textbf{约束}：
\begin{align}
F_1^T (V_A - V_B) &\geq 0 \quad \text{（接触1）} \\
F_2^T (V_B - V_C) &\geq 0 \quad \text{（接触2）}
\end{align}

\textbf{可行twist集合}：
\begin{itemize}
\item 复合twist空间：$(V_A, V_B, V_C) \in \mathbb{R}^{18}$
\item 两个约束在18维空间中定义了两个半空间
\item 可行集合是两个半空间的交集：一个以原点为根部的凸锥
\end{itemize}

\section{维度分析}

\subsection{复合twist空间的维度}

对于$m$个物体：
\begin{itemize}
\item \textbf{3D情况}：每个物体6维，总共$6m$维
\item \textbf{平面情况}：每个物体3维，总共$3m$维
\end{itemize}

\subsection{约束的维度}

每个成对接触约束：
\begin{itemize}
\item 在$6m$维（或$3m$维）复合twist空间中
\item 定义了一个$(6m-1)$维（或$(3m-1)$维）的半空间
\item 边界是$(6m-1)$维（或$(3m-1)$维）的超平面
\end{itemize}

\subsection{可行集合的维度}

\begin{itemize}
\item 如果有$k$个独立的接触约束
\item 可行集合的维度至少是$6m - k$（或$3m - k$）
\item 如果所有约束都主动，维度可能更小
\end{itemize}

\section{齐次 vs 非齐次约束}

\subsection{齐次约束（所有物体非受控）}

如果所有物体都非受控：
\begin{itemize}
\item 约束形式：$F^T (V_i - V_j) \geq 0$
\item 这是齐次约束（超平面通过原点）
\item 可行集合：以原点为根部的多面体凸锥
\item 原点总是可行的（所有物体静止）
\end{itemize}

\subsection{非齐次约束（有受控物体）}

如果有受控物体：
\begin{itemize}
\item 约束形式：$F^T (V_i - V_j) \geq 0$，其中$V_j$是受控的
\item 可以写成：$F^T V_i \geq F^T V_j$（非齐次）
\item 可行集合：一般的多面体（可能是有界的）
\item 原点可能不可行（如果受控物体在运动）
\end{itemize}

\section{关键概念总结}

\subsection{从单个物体到多个物体}

\begin{center}
\begin{tabular}{|c|c|c|}
\hline
\textbf{方面} & \textbf{单个物体（12.1.3）} & \textbf{多个物体（12.1.4）} \\
\hline
考虑对象 & 物体$A$的twist $V_A$ & 所有物体的twist $(V_1, \ldots, V_m)$ \\
\hline
空间维度 & 6维（$V_A$空间） & $6m$维（复合twist空间） \\
\hline
约束形式 & $F_i^T (V_A - V_{j(i)}) \geq 0$ & $F^T (V_i - V_j) \geq 0$ \\
\hline
可行集合 & $V_A$空间中的多面体 & 复合空间中的多面体 \\
\hline
\end{tabular}
\end{center}

\subsection{受控 vs 非受控}

\begin{center}
\begin{tabular}{|c|c|c|}
\hline
\textbf{特征} & \textbf{所有物体非受控} & \textbf{有受控物体} \\
\hline
约束类型 & 齐次 & 非齐次 \\
\hline
可行集合 & 以原点为根部的锥 & 一般多面体（可能封闭） \\
\hline
原点 & 总是可行 & 可能不可行 \\
\hline
\end{tabular}
\end{center}

\section{物理直觉}

\subsection{为什么需要复合twist空间？}

\begin{itemize}
\item 每个物体都有自己的运动
\item 接触约束涉及\textbf{两个物体}的相对运动
\item 需要在\textbf{所有物体的联合空间}中考虑约束
\item 就像在"大房间"中考虑所有物体的位置，而不是单独考虑每个物体
\end{itemize}

\subsection{为什么受控物体会改变约束性质？}

\begin{itemize}
\item 受控物体的运动是"固定的"
\item 约束变成了对非受控物体的"绝对限制"
\item 就像"墙"在移动，而不是静止的
\item 可行区域会"移动"，不再是固定的锥
\end{itemize}

\section{实际应用}

\subsection{抓取操作}

在机器人抓取中：
\begin{itemize}
\item 手指是受控物体（运动由机器人控制）
\item 被抓取物体是非受控物体（运动由接触约束决定）
\item 需要分析被抓取物体在给定手指运动下的可能运动
\end{itemize}

\subsection{装配分析}

在装配分析中：
\begin{itemize}
\item 所有物体可能都非受控
\item 需要分析整个装配的可行运动
\item 判断装配是否稳定（形状闭合）
\end{itemize}

\section{总结}

12.1.4章节的核心思想：

\begin{enumerate}
\item \textbf{推广到多个物体}：从单个物体扩展到多个物体同时接触的情况

\item \textbf{复合twist空间}：需要在所有物体的联合twist空间中考虑约束

\item \textbf{成对接触约束}：每个接触约束涉及两个物体的相对运动

\item \textbf{受控 vs 非受控}：
  \begin{itemize}
  \item 所有物体非受控：齐次约束，可行集合是锥
  \item 有受控物体：非齐次约束，可行集合是一般多面体
  \end{itemize}

\item \textbf{接触模式}：整个装配的接触模式是所有接触标签的连接
\end{enumerate}

理解这些概念对于分析复杂的多物体接触系统至关重要。

\end{document}

